% ============================================================
% Преамбула для PDF-версии лекций (XeLaTeX)
% ============================================================

% Русский язык
\usepackage{polyglossia}
\setdefaultlanguage{russian}
\setotherlanguage{english}

% Шрифты
\usepackage{fontspec}
\setmainfont{PT Serif}[
  BoldFont      = PT Serif Bold,
  ItalicFont    = PT Serif Italic,
  BoldItalicFont= PT Serif Bold Italic
]
\setsansfont{PT Sans}
\setmonofont{DejaVu Sans Mono}[Scale=0.9]

% Математика
\usepackage{amsmath, amsthm, amssymb, mathtools}

% Теоремы на русском
\newtheorem{theorem}{Теорема}[section]
\newtheorem{lemma}[theorem]{Лемма}
\newtheorem{corollary}[theorem]{Следствие}
\theoremstyle{definition}
\newtheorem{definition}[theorem]{Определение}
\newtheorem{example}[theorem]{Пример}
\theoremstyle{remark}
\newtheorem*{remark}{Замечание}

% Красивые заголовки
\usepackage{titlesec}
\titleformat{\section}{\Large\bfseries\sffamily}{{\thesection}}{1em}{}
\titleformat{\subsection}{\large\bfseries\sffamily}{{\thesubsection}}{1em}{}

% Колонтитулы
\usepackage{fancyhdr}
\pagestyle{fancy}
\fancyhf{}
\fancyhead[L]{\small\leftmark}
\fancyhead[R]{\small Математический анализ}
\fancyfoot[C]{\thepage}

% Цвета и гиперссылки
\usepackage{xcolor}
\definecolor{mathblue}{HTML}{2563eb}
\usepackage{hyperref}
\hypersetup{
  colorlinks=true,
  linkcolor=mathblue,
  urlcolor=mathblue,
  citecolor=mathblue
}

% Красивые блоки для доказательств и замечаний
\usepackage{tcolorbox}
\tcbuselibrary{theorems, breakable, skins}

\tcolorboxenvironment{theorem}{
  colback=blue!5, colframe=blue!50!black,
  boxrule=0.5pt, breakable, enhanced
}
\tcolorboxenvironment{definition}{
  colback=green!5, colframe=green!50!black,
  boxrule=0.5pt, breakable, enhanced
}
